%=============================================================================!
% 10.1  A BOX REPEATED IN EVERY DIRECTION                                     !
%=============================================================================!
\section{A box repeated in every direction}
\label{sec:md-periodic}

A sample small enough to simulate is mostly surface. This section counts
how much, removes the surface by repeating a box of atoms in every
direction, and sets up the coordinates that describe such a box.

\subsection*{Most of a small sample is surface}

We place $N = n^3$ atoms on a cubic grid, $n$ along each edge. The atoms that
are not on a face form a smaller cube of $(n - 2)^3$, so the fraction on
the surface is
\begin{align*}
   1 - \frac{(n - 2)^3}{n^3} &= 1 - \Bigl(1 - \frac2n\Bigr)^3\\
   &= \frac6n - \frac{12}{n^2} + \frac8{n^3}
   && \text{expanding $(1 - x)^3 = 1 - 3x + 3x^2 - x^3$ with $x = 2/n$,}
\end{align*}
close to $6/n$ when $n$ is large. A thousand atoms ($n = 10$) have 49\%
of them on the surface, about ten thousand ($n = 22$) 25\%, and a million ($n =
100$) still 5.9\%. A surface atom has fewer neighbours than one inside,
and a sample this small would measure the properties of its surface as
much as those of its bulk.

\subsection*{The periodic box}

The surface is removed by repeating the box. The atoms sit in a cell,
and the cell and its contents are copied without end in every direction,
so that space is filled with identical cells. An atom that leaves through
one face of the cell comes in through the opposite face, as its copy from
the neighbouring cell enters, and every atom has neighbours on all sides;
a ship in an old arcade game that leaves the screen at the right edge and
reappears at the left lives in such a box, in two dimensions. Each atom
feels the others and their copies. The price is an imposed periodicity:
the sample repeats with the period of the cell, no wave longer than the
cell, such as a long sound wave, can occur, and for a short-ranged
potential the cell must be large enough that an atom does not feel its
own copies (\cref{sec:md-image}); the Coulomb energy of
\cref{sec:md-ewald} includes them by design.

The cell is a parallelepiped, a box whose six faces are parallelograms,
spanned by three lattice vectors $\vb{a}$,
$\vb{b}$ and $\vb{c}$, which we place as the columns of the cell matrix
\begin{equation}
   \vb{h} = \begin{pmatrix} \vb{a} & \vb{b} & \vb{c} \end{pmatrix} .
   \label{eq:md-cell}
\end{equation}
For a cubic cell of side $L$, $\vb{a} = (L, 0, 0)$, $\vb{b} = (0, L, 0)$,
$\vb{c} = (0, 0, L)$ and $\vb{h}$ is $L$ times the identity. Any position
can be written as $\vb{r} = s_1\vb{a} + s_2\vb{b} + s_3\vb{c}$, which is
$\vb{h}$ times the vector $\vb{s} = (s_1, s_2, s_3)$: row $x$ of $\vb{h}$
is $(a_x, b_x, c_x)$, so the $x$ component of $\vb{h}\vb{s}$ is $a_xs_1 +
b_xs_2 + c_xs_3$, the $x$ component of $s_1\vb{a} + s_2\vb{b} +
s_3\vb{c}$, and likewise for $y$ and $z$. So
\begin{equation}
   \vb{r} = \vb{h}\vb{s} , \qquad \vb{s} = \vb{h}^{-1}\vb{r} ,
   \label{eq:md-fractional}
\end{equation}
with the inverse matrix of \cref{sec:ro-remove}. The inverse exists: a
matrix without one turns some $\vb{s} \neq \vb{0}$ into $\vb{0}$, and
$s_1\vb{a} + s_2\vb{b} + s_3\vb{c} = \vb{0}$ would put the three lattice
vectors in one plane, a flat cell with no volume. The numbers $s_1$, $s_2$
and $s_3$ are the fractional coordinates of the position: the point lies
inside the cell when each is between 0 and 1, and its copies are at
$\vb{h}(\vb{s} + \vb{n})$, with $\vb{n}$ any vector of three whole
numbers, positive, negative or zero.

The volume of the cell is its base times its height. The face spanned by
$\vb{b}$ and $\vb{c}$ is a parallelogram of area $\lvert\vb{b}\times\vb{c}
\rvert$ (\cref{sec:ro-cross}), and $\vb{b}\times\vb{c}$ is at right
angles to that face. The height of the cell above it is the size of the
part of $\vb{a}$ along that direction, the scalar product of $\vb{a}$
with the vector of length one $(\vb{b}\times\vb{c})/\lvert\vb{b}\times
\vb{c}\rvert$ (\cref{sec:mo-plane}), taken without its sign, since
$\vb{a}$ may lie on either side of the face. Then
\begin{align}
   V &= \lvert\vb{b}\times\vb{c}\rvert\,
        \frac{\lvert\vb{a}\cdot(\vb{b}\times\vb{c})\rvert}
             {\lvert\vb{b}\times\vb{c}\rvert}
     && \text{base times height}\notag\\
     &= \lvert\vb{a}\cdot(\vb{b}\times\vb{c})\rvert ,
     && \text{cancelling $\lvert\vb{b}\times\vb{c}\rvert$}
   \label{eq:md-volume}
\end{align}
the scalar triple product of \cref{sec:ro-tensor}. For the cubic cell, $V =
L^3$.

\subsection*{Wrapping and unwrapping}

To put a position back into the cell, we take its fractional coordinates
and subtract from each its floor,
\begin{equation*}
   \vb{s} \to \vb{s} - \lfloor\vb{s}\rfloor ,
\end{equation*}
where the floor $\lfloor x\rfloor$ of a number $x$ is the largest whole
number not above it, and $\lfloor\vb{s}\rfloor$ takes the floor of each
coordinate: $\lfloor 2.3\rfloor = 2$ and $\lfloor -0.4\rfloor = -1$, so
2.3 becomes 0.3 and $-0.4$ becomes 0.6.
This is wrapping. The function \texttt{wrap} of \texttt{mdlab.cell} does
it for $N$ atoms:

\begin{code}
def wrap(positions: ArrayLike, cell: ArrayLike) -> NDArray:
    s = to_fractional(positions, cell)
    return to_cartesian(s - np.floor(s), cell)
\end{code}

\noindent where \texttt{to\_fractional} solves $\vb{h}\vb{s} = \vb{r}$
for each atom and \texttt{to\_cartesian} multiplies back.

The integration itself does not need wrapped positions: the forces depend
only on separations, and \cref{sec:md-image} reduces each separation to
its nearest copy whichever copy of each atom the code holds. \texttt{mdlab} therefore follows
each atom's path without wrapping, so that its displacement over a run can
be read off directly, and wraps the positions only when it writes them to
a file, as files of periodic systems usually hold them. To recover a
continuous path from wrapped frames, each frame's position is replaced by
its copy nearest to the previous frame's, which is right as long as no
atom moves more than half the narrowest width of the cell between frames
(\cref{sec:md-image}): a longer move is mistaken for a shorter one the
other way, as too few frames per period hid an oscillation in
\cref{sec:mo-atoms}.

\subsection*{Momentum in a periodic box}

Moving every atom by the same displacement moves each of its copies with
it and changes no separation, so the forces still add to zero and the
total momentum is conserved, as \cref{sec:ha-noether} found, whether a
pair interacts through its nearest copy (\cref{sec:md-image}) or through
all of them (\cref{sec:md-ewald}). Turning the atoms is not
a symmetry of the box, and the angular momentum is not conserved. In the
run of 500 atoms of \cref{sec:md-loop} the largest component of the total
momentum after 2000 steps is $\num{3.8e-15}$\,amu\,\AA/fs, the rounding of
the arithmetic.

\begin{keyidea}
A cell repeated in every direction removes the surface that would
dominate any sample small enough to simulate. The cell is described by
the matrix $\vb{h}$ of its lattice vectors; positions have fractional
coordinates $\vb{s} = \vb{h}^{-1}\vb{r}$, the cell has volume $\lvert
\vb{a}\cdot(\vb{b}\times\vb{c})\rvert$, and wrapping subtracts from
each fractional coordinate its floor.
\end{keyidea}

\notebook{10}{wrap atoms into a cubic and a hexagonal cell}

\begin{selfcheck}
In the cell with $\vb{a} = (4, 0, 0)$, $\vb{b} = (0, 5, 0)$ and $\vb{c} =
(0, 0, 6)$\,\AA, what are the fractional coordinates of $\vb{r} = (9,
-1, 3)$\,\AA, and where does wrapping put it?
\end{selfcheck}
